\chapter{Continuation on the Line and Banach -Szidon
  Sequences}\label{chap18} %cha 18

We\pageoriginale have seen in the last lecture that if $f \in \mathscr{C}_\Lambda
(I), | I | > 2 \pi \triangle$ and if $f$ belongs to a class of
$C^\infty$- functions, then $f$ can be continued to a function $\in
\mathscr{C}_\Lambda$. The answer without the assumption of
differentiability is not known completely. But conditions on $\Lambda$
can be given to have the answer in particular cases. Suppose $\Lambda$
is real, symmetric and very lacunary in the sense that $\lambda_{n +
 1}|. \lambda_n >> 3$. Then we shall prove that every function $f \in
\mathscr{C}_\Lambda (I)$ is continuable into $\mathscr{C}_\Lambda$. We
take a trigonometric polynomial $P(x) = \sum\limits_{n_o} r_n \cos
(\lambda_n x + \varphi_n), \big | r_n \big | \ge 0$. 

For $\lambda_n$ large enough the set $E_n$ is defined as the set of
points $x$ satisfying the following conditions: 
$$
\pi/2 - \alpha + 2 k \pi < \lambda_n x+ \varphi_n < \pi/2 + \alpha + 2
k \pi (k =. . - 1, 0, 1, \ldots) 
$$

The connected components of $E_n$ are intervals of length
$\dfrac{\pi - 2 \alpha}{\lambda_n}$. 

\begin{figure}[H]
 \centerline{\includegraphics{vol15-figures/fig15-20a.eps}}
\end{figure}

If $n_o$ is large enough, $E_{n_o} $ has at least one connected
component entirely in $I$, say $I_{n_o}$. Assuming $\dfrac{\lambda_
{n+1}}{\lambda_n} \ge 1 + 2 \dfrac{\pi + 2 \alpha}{\pi - 2 \alpha}$,
there is at least one connected component of $E_{n + 1}$, say $I_{n +1
}$, inside a connected component $I_n $ of $E_n$ (see fig.) ; we
define, in that way, $I_{n_o+1}$. $I_{n_o + 2}, \ldots, $ Take $X \in
\cap I_n $; thus 
\begin{equation}
 P(X) > \sin \alpha \sum^N_{n_o} r_n, \tag{1}\label{chap18:eq1}
\end{equation}

Let $\Lambda' = \Lambda - \big \{ \pm \lambda_1, \pm \lambda_2, \ldots
, \pm \lambda_{n_o-1} \big \}$, $\big | f \big |_I = \max\limits_{x
 \in I} \big | f(x) \big |$ and $ || f || = \sum r_n, $ if $f =
\sum\limits_{n_o}^\infty r_n \cos (\lambda_n x + \varphi_n), \sum r_n
= \sum \big | r_n \big | < \infty$; they\pageoriginale are respectively norms in two
Banach spaces, say $\mathscr{C} (I)$ and $A_{\Lambda'}$. (1) proves
the equivalence of these norms for the polynomials in
$\mathscr{C}_{\Lambda'} (I)$: $\sin \alpha || P || < | P |_I \le || P
||$. As the polynomials $P$ form a dense subset in
$\mathscr{C}_{\Lambda'}(I)$ and in $A_{\Lambda'}$, we can identify
$\mathscr{C}_\Lambda' (I)$ and $A_\Lambda'$: every $f \in
\mathscr{C}_{\Lambda'}(I)$ can be expressed as an element of
$A_{\Lambda'}$. The same result holds by adding a finite number of
terms to $\Lambda'$ (because 
$$
f \in \mathscr{C}_\Lambda (I) \Leftrightarrow f=f_1 + f_2, f_1 = \sum_1^{n_0-1}
r_n \cos (\lambda_n x + \varphi_n), f \in \mathscr{C}_{\Lambda'} (I)
). 
$$

Thus we have prove the following proposition:

\begin{prop*} % pro
 Suppose $\Lambda = \big \{ \pm \lambda_n \big\}$ is a real symmetric
 sequence which is lacunary in the sense that $\lambda_{n + 1}/
 \lambda_n >> 3$. Then every function $f \in \mathscr{C} (I)$ is the
 sum of an absolutely convergent Fourier-series and $f$ is
 continuable to $\mathscr{C}_\Lambda$. 
\end{prop*}
 
A result of this type was considered by Szidon (Zygmund 1, chap. VI
\S 6.4) who has proved that a bounded periodic function whose
spectrum $\Lambda$ is lacunary in the sense of Hadamard,
viz. $\lambda_{n + 1} / \lambda_n >> 1$, has an absolutely convergent
Fourier series. 

\begin{defi*}
 A sequence $\Lambda$ is said to be a Szidon sequence if for every
 interval $I, f \in \mathscr{C} (I) \Rightarrow f = \sum a (\lambda)
 e^{i \lambda x}$ and $\sum \big | a (\lambda) \big | < \infty$. 
\end{defi*}

We shall make a brief study of Szidon sequences and give alternate
definitions of them by a simple lemma on Banach spaces. 

Let\pageoriginale $\mathscr{E}_1$ and $\mathscr{E}_2$ be two Banach spaces with
norms $|| ~ ||_1$ and $|| ~ ||_2$ and let $\mathscr{E}_1$ be an
algebraic subspace of $\mathscr{E}_2$. Moreover, for every $x \in
\mathscr{E}_1$, let $|| x ||_1 \ge || x ||_2$. Then $\mathscr{E}'_2$
is imbedded in $\mathscr{E}'_1$ by $\langle x, x' \rangle_1 = \langle
x, x' \rangle_2 \mathscr{E}'_2 \subset \mathscr{E}'$. 

By a theorem of Banach, if $\mathscr{E}_1 = \mathscr{E}_2 $ then the
norms are equivalent and $\mathscr{E}'_1 =
\mathscr{E}'_2$. Conversely, if $\mathscr{E}'_1 = \mathscr{E}'_2$,
then is it true that $\mathscr{E}_1 = \mathscr{E}_2$ ? But $|| x ||_1 =
\sup\limits_{|| x' || \le 1} \big | \langle x, x' \rangle |, || x ||_2
= \sup_{|| x ' || \le 1} \big | \langle x, x' \rangle \big |$. Since
$\mathscr{E}'_1 = \mathscr{E}'_2$ and since $|| x' ||_2 \ge K^{-1} ||
x' ||_1, K > 0$, we have $|| x ||_1 < K \sup\limits_{|| x' ||} \le
\big | \langle x, x' \rangle \big | = K || x ||_2$. This means that
$\mathscr{E}_1$ is closed in $\mathscr{E}_2$. In order that
$\mathscr{E}_1 = \mathscr{E}_2$ it is necessary and sufficient that
$\mathscr{E}_1$ is dense in $\mathscr{E}_2$. That proves 
\begin{lemma*} % lem 
  Let $\mathscr{E}_1 $ and $\mathscr{E}_2$ be two Banach spaces with
  norms $|| ~ ||_1$ and $|| ~ ||_2$; suppose $\mathscr{E}_1$ is an
  algebraic subspace of $\mathscr{E}_2, \mathscr{E}_1$ is dense in
  $\mathscr{E}_2$ and $|| x ||_1 \le || x ||_2$ for every $x \in
  \mathscr{E}'_1$. If $\mathscr{E}'_1 = \mathscr{E}'_2$, then
  $\mathscr{E}_1 = \mathscr{E}_2$.
\end{lemma*}

\begin{defi*} % def
 We define a sequence $\Lambda$ to be a {\em Banach sequence } if for
 every interval $I$ and every sequence $b (\lambda) \to 0$ as
 $\lambda (\in \Lambda) \to \infty$ there exists a function $f \in
 L^1 (I)$ such that $b (\lambda) = \int_I f e^{- i \lambda x} dx \,(~
 \lambda \in \Lambda)$. 
\end{defi*}

We prove a theorem relating Banach sequences and Szidon sequen\-ces, by
applying the lemma. 

Let $\mathscr{E}_2$ be the space of sequences $\big \{ a( \lambda)
\big \}$ converging to zero endowed with the norm, $|| \{ a (\lambda
)\} ||_2 = \sup\limits_{\lambda} \big | a (\lambda) \big |$. 

Let $\mathscr{E}_1$ be the space of sequences $\{b (\lambda) \}$ with
$b (\lambda) = \int_I f e^{-i \lambda x} dx$, where $\int_I \big | f
(x) \big | dx < \infty$. The norm in $\mathscr{E}_1$ is taken as $||
\{b (\lambda) \} ||_1 = \inf \int_I | f |$ where the infimum is taken
over all $f$ for which $b(\lambda) = \int_I e^{-i \lambda x} f (x)
dx$. $\mathscr{E}_1$ is isomorphic to $L^1 (I) / H$, where $H$ is the
subspace of $L^1 (I)$ orthogonal to $\{e^{i \lambda x} \}$. 

$\mathscr{E}'_2$ is the space of absolutely convergent sequences $\{ C
(\lambda) \}$, with the bilinear form on $\mathscr{E}_2 x
\mathscr{E}'_2$ as $\langle a(\lambda), C (\lambda) \rangle = \sum a
(\lambda) C (\lambda)$. The norm in $\mathscr{E}'_2$ is given by $|| C
(\lambda) ||_2 = \sum C(\lambda). \mathscr{E}'_2$ is isomorphic to
$A_\Lambda$. 

$\mathscr{E}'_1$ is isomorphic to the subspace $\mathscr{C}_\Lambda
(I)$ of $L^\infty (I)$ spanned by $\big \{ e^{i \lambda x} \big\}$,
with the norm of $\{d (\lambda)\} = || \varphi ||_{L^\infty (I)}$
where $\varphi \sim \sum d (\lambda) e^{i \lambda x}$. The bilinear
form giving the duality is : $\langle b (\lambda), d (\lambda) \rangle
= \int_I f \varphi$, where $b(\lambda) = \int_I f(x) e^{-i \lambda x }
dx$. 

Now\pageoriginale the sequence $\Lambda$ is a Banach sequence if and only if
$\mathscr{E}_1 = \mathscr{E}_2$, and a Szidon sequence if and only if
$\mathscr{E}'_1 = \mathscr{E}'_2$. A Banach sequence $\Lambda$ is a
Szidon sequence. In order to see that a Szidon s equence be a Banach
sequence, it is sufficient to prove that $\mathscr{E}_1$ is dense in
$\mathscr{E}_2$. This is so since every linear functional on
$\mathscr{E}_2$ which is zero on $\mathscr{E}_1$ is zero. We can apply
again our lemma, taking $(\mathscr{E}'_1, \mathscr{E}'_2)$ instead of
$(\mathscr{E}'_2, \mathscr{E}'_1)$. 

$\mathscr{E}''_2$ is the space of bounded sequences $\big \{ e
(\lambda) \big\}$ with $\langle C(\lambda), e (\lambda) \rangle = \sum
c (\lambda) e(\lambda)$ and the norm $|| e (\lambda) || = \sup \big |
e(\lambda) \big |. \mathscr{E}''_1$ is isomorphic to the quotient space
$M(I) /K$ where $M(I) $ is the space of measures with support in $I$
and $K$ is the subspace of it orthogonal to $\big \{e^{i \lambda x}
\big \}$. $\big \{ f (\lambda) \big \} \in \mathscr{E}''_1$ when $f
(\lambda) = \int_I e^{i \lambda x} d \mu, d \mu \in M(I) /K$ and $|| b
(\lambda) || = \inf\limits_{\mu} \int | d \mu |$. When $\mathscr{E}_1
= \mathscr{E}_2, \mathscr{E}''_1 = \mathscr{E}''_2$. Moreover
$\mathscr{E}'_2$ is dense 
in $\mathscr{E}'_1$ (because obviously $A_\Lambda$ is dense in
$\mathscr{C}_ \Lambda (I))$; then $\mathscr{E}'_1 = \mathscr{E}'_2
\Rightarrow \mathscr{E}''_1 = \mathscr{E}''_2$. We get the following
theorem: 
\begin{theorem*}
 The definitions of Szidon sequences and Banach sequences are
 equivalent, and equivalent to the following: for every integral $I$
 and for every bounded sequence $\big \{e (\lambda) \big \} ( \lambda
 \in \Lambda)$, there exists a measure $d \mu$ with support in $I$,
 such that $e(\lambda) = \int_I e^{-i \lambda x} d \mu (x) (\lambda
 \in \Lambda)$. 
\end{theorem*}

We have given sufficient conditions for a sequence $\{\pm \lambda_n\}$
to be a Banach-Szidon sequence. Actually, Zygmund has shown (Zygmund
$2$) that the condition $\dfrac{\lambda_{n+1}}{\lambda_n} > > 1$ is
sufficient. Necessary conditions involve arithmetical properties of
$\lambda_n$, as it appears from the following proposition (proved in
(Kahane $2$) ~ ). 

Let\pageoriginale $\xi_1, \ldots, \xi_p$ be a real numbers. Consider $n_1 \xi_1 +
\cdots + n_p \xi_p, n_j$ being integers satisfying the inequality $|
n_1 | + \cdots + |n_p | \le s$. If $\Lambda$ is a Banach-Szidon
sequence, there exists a constant $A = A(\lambda)$ such that among the
set of $n_1 \xi_1 + \cdots + n_p \xi_p$ there are only $[A_p \log
 (1+s)]$ of them in $\Lambda$. 

\begin{examples*}
 Take $p = 1$. We cannot have more than $0 (\log s)$ points of
 $\Lambda$ in an arithmetical progression containing $s$ terms. 
\end{examples*}

Let $s = 2$, and $p = 1, 2, \ldots$ Whatever be the sequence $\{
\xi_j\}$, we cannot have in $\Lambda$ more than $0(p)$ points among the
$p^2$ points $\xi_j + \xi_k$, $| j | \le p, | k | \le p$. 

The study of Banach-Szidon sequences is interesting in itself, but can
give only a partial answer to the following problem: 

\noindent
\textbf{Problem.} Give a relation, between $\Lambda$ and $I$, such
that every $f \in \mathscr{C}_\Lambda (I)$ is the restriction on $I$
of a function $\in \mathscr{C}_\Lambda$. 

This problem becomes easier and has pretty good solutions if we
replace $\mathscr{C}$ by the space of functions which belong to $L^2$
on every interval, with the topology of the convergence in $L^2$ on
every interval ( (Paley - Wiener), (Kahane $2$) ). 
